\scanpage{41}

C'est là une condition \inserted{de décroissance des coefficients} arithmétique assez biscornue. W.M. la remplacent par une condition plus forte, et plus maniable, savoir que
\[
|a_n|=O(1)\rho^n
\qquad
\begin{gathered}
\text{pour un }\rho>1\\
\text{convenable (dépendant de la suite }(a_n)\text{)}
\end{gathered}
\]
Cela signifie qu'au lieu de regarder l'anneau $W\{t\}$ des séries convergentes sur le disque unité \inserted{de $\overline K$ ($K$ corps des fractions de $W$)}, ils regardent l'anneau \inserted{$W^{\unsure{\dagger|+}}\{t\}$} des séries convergentes sur un disque de rayon $\rho>1$ \deleted{(\unsure{de rayon arbitraire}} un peu plus grand (dépendant de la fonction envisagée), On définit de même $W^{\unsure{\dagger|+}}\{t_1,\ldots,t_n\}$, et les \deleted{\unsure{anneaux}} W-M-algèbres, \deleted{de W.M} qui sont les quotients de tels anneaux par des idéaux \deleted{quelconques} fermés quelconques \deleted{(\unsure{fermés}} (les idéaux sont-ils

\margin{40}
